\section{Latar Belakang}
Compiler adalah perangkat lunak yang menerjemahkan source code dari bahasa 
tingkat tinggi menjadi representasi yang dapat dieksekusi oleh mesin. Keberadaan 
compiler penting karena memungkinkan programmer menulis program dengan struktur dan 
ekspresi yang mendekati cara manusia berpikir, namun tetap menghasilkan instruksi 
yang dapat dipahami oleh komputer. Dalam praktik rekayasa perangkat lunak, compiler 
tidak hanya berfungsi sebagai penerjemah, tetapi juga sebagai penjaga konsistensi 
aturan bahasa, sehingga program yang dihasilkan bersifat benar dan dapat dijalankan.

Proses kompilasi terdiri dari beberapa tahap yang saling bergantung. Tahap lexical 
analysis memecah karakter menjadi token; syntax analysis memeriksa apakah urutan 
token membentuk struktur yang sah; semantic analysis memvalidasi makna, tipe data, 
dan aturan konteks; sedangkan code generation mengubah representasi internal menjadi 
bentuk target. Setiap tahap memiliki peran yang kritis karena kesalahan pada tahap 
awal akan menghambat atau merusak proses pada tahap berikutnya. Oleh sebab itu, 
pembangunan compiler yang benar menuntut setiap tahap dikerjakan secara sistematis.

Pada mata kuliah IF2224 Teori Bahasa Formal dan Otomata, tugas besar Arion Compiler 
dirancang untuk menghubungkan teori dengan implementasi nyata. Bahasa Arion yang 
bersifat Pascal-like dipilih agar mahasiswa dapat menerapkan konsep grammar formal, 
otomata, serta teknik parsing pada bahasa yang memiliki struktur jelas. Pengerjaan 
bertahap melalui milestone mendorong pemahaman yang mendalam atas keterkaitan antar 
komponen compiler, sekaligus memberi ruang untuk menguji hasil secara terpisah pada 
setiap fase.

Milestone 1 telah menyelesaikan pembangunan lexer yang menghasilkan daftar token 
dari source code Arion. Token inilah yang menjadi masukan utama bagi Milestone 2, 
yaitu pembangunan parser atau syntax analyzer. Parser bertugas memeriksa struktur 
sintaksis program dan memastikan token-token tersebut mengikuti grammar yang telah 
ditentukan, serta membangun parse tree sebagai representasi hierarkisnya. Dengan 
demikian, Milestone 2 menjadi jembatan antara tokenisasi dan analisis semantik yang 
akan dilaksanakan pada milestone berikutnya.

Pembangunan parser pada milestone ini sangat terkait dengan teori Context-Free 
Grammar (CFG) dan teknik Recursive Descent. CFG mendefinisikan struktur bahasa dalam 
bentuk produksi non-terminal ke terminal dan non-terminal, sedangkan Recursive Descent 
merealisasikan setiap non-terminal sebagai fungsi rekursif dalam kode C/C++. Keduanya 
memberikan pendekatan yang langsung dan transparan dalam mengimplementasikan parser, 
sehingga proses analisis sintaks dapat ditelusuri sesuai dengan kaidah formal yang dipelajari.

\section{Rumusan Masalah}
Penyusunan parser untuk Arion Compiler memunculkan kebutuhan untuk merumuskan pertanyaan 
penelitian yang menjembatani teori dan implementasi. Laporan ini berangkat dari pertanyaan 
utama mengenai bagaimana membangun parser yang mampu menganalisis struktur sintaksis dari 
daftar token hasil lexer sehingga sesuai dengan grammar bahasa Arion. Pertanyaan ini mencakup 
bagaimana token-token disusun ulang menjadi konstruksi program yang valid dan bagaimana aturan 
grammar ditegakkan secara konsisten.

Secara lebih teknis, tantangan berikutnya adalah bagaimana mengimplementasikan algoritma 
Recursive Descent agar setiap produksi grammar dapat direpresentasikan dengan jelas 
dalam kode C/C++. Implementasi ini harus memperhatikan urutan produksi, pemilihan 
cabang grammar berdasarkan token lookahead, serta pemetaan non-terminal menjadi 
fungsi-fungsi parser yang terstruktur dan mudah dilacak.

Selain itu, laporan ini juga mengkaji bagaimana membangun dan menampilkan parse 
tree yang menggambarkan hierarki struktur program Arion. Parse tree tidak hanya 
berfungsi sebagai keluaran parser, tetapi juga sebagai jembatan menuju tahap semantic 
analysis pada milestone selanjutnya. Oleh karena itu, representasi parse tree harus 
akurat, lengkap, dan konsisten terhadap grammar.

Aspek lain yang tidak kalah penting adalah bagaimana parser menangani kesalahan 
sintaksis secara informatif. Parser harus mampu mendeteksi token yang tidak sesuai, 
menunjukkan lokasi kesalahan, serta memberikan pesan yang jelas agar pengguna dapat 
memperbaiki program. Pertanyaan mengenai strategi error handling ini menjadi bagian 
penting dalam memastikan parser bersifat robust dan dapat digunakan pada pengujian nyata.

\section{Tujuan}
Tujuan utama Milestone 2 adalah membangun komponen parser Arion Compiler yang dapat 
menganalisis daftar token dari Milestone 1 menggunakan algoritma Recursive Descent. 
Parser yang dibangun harus mampu memvalidasi struktur program secara sintaktis dan 
menghasilkan parse tree sebagai keluaran utama. Pencapaian ini menjadi fondasi agar 
pipeline kompilasi dapat berlanjut ke tahap semantic analysis pada milestone berikutnya.

Tujuan berikutnya adalah merumuskan dan mengimplementasikan grammar bahasa Arion secara 
terstruktur di dalam program C/C++. Grammar tersebut harus mencakup seluruh konstruksi 
yang ditentukan dalam spesifikasi Milestone 2, sehingga parser memiliki cakupan yang 
tepat terhadap fitur bahasa Arion. Dengan implementasi yang terstruktur, grammar akan 
mudah ditelusuri dan dijadikan acuan saat melakukan pengujian.

Laporan ini juga bertujuan menghasilkan parse tree yang merepresentasikan struktur 
hierarkis program Arion secara lengkap. Parse tree harus dapat dicetak ke terminal 
dan disimpan ke file berformat .txt sesuai spesifikasi, sehingga hasil parsing 
dapat diverifikasi secara manual maupun digunakan kembali pada tahap berikutnya.

Tujuan lain yang penting adalah menyediakan mekanisme error handling pada level parser. 
Parser harus mendeteksi syntax error, menandai posisi kesalahan, dan menampilkan 
pesan yang informatif agar proses debugging lebih efisien. Selain itu, tujuan 
akhir dari milestone ini adalah memastikan bahwa keluaran parse tree dapat digunakan 
sebagai masukan langsung untuk tahap semantic analysis yang akan datang.

\section{Batasan Masalah}
Pengerjaan Milestone 2 dibatasi oleh spesifikasi bahasa Arion yang ditetapkan pada 
dokumen resmi hingga Revisi 5 tertanggal 04-05-2026. Parser yang dibangun hanya 
mengakomodasi grammar dan aturan bahasa Arion pada milestone ini, sehingga tidak 
mencakup bahasa lain maupun fitur di luar spesifikasi. Batasan ini menjaga fokus 
implementasi agar sesuai dengan kebutuhan tugas besar.

Algoritma parsing yang digunakan terbatas pada Recursive Descent. Pendekatan ini 
dipilih karena kesederhanaan implementasi dan kesesuaiannya dengan grammar yang 
tersedia, sehingga algoritma lain seperti LR, LALR, atau Earley tidak digunakan. 
Konsekuensinya, grammar harus disesuaikan agar kompatibel dengan parsing top-down.

Keluaran utama parser adalah parse tree, bukan abstract syntax tree (AST). Parse 
tree memuat seluruh simbol terminal dan non-terminal sesuai grammar, sementara 
AST akan menjadi fokus pada milestone berikutnya. Selain itu, implementasi dilakukan 
dalam bahasa C/C++, melanjutkan struktur kode Milestone 1, agar integrasi antara 
lexer dan parser tetap konsisten.

Masukan parser dibatasi pada file daftar token hasil lexer Milestone 1, bukan langsung 
dari source code Arion. Meskipun secara pipeline lexer dan parser dijalankan berurutan, 
parser tetap menerima format token stream sesuai spesifikasi. Terakhir, parser pada 
milestone ini hanya melakukan pengecekan sintaksis dan belum mencakup analisis semantik 
seperti pengecekan tipe data atau scope variabel.
